\chapter{Architektura środowiska laboratoryjnego i metodyka badawcza}\label{ch:architektura_laboratorium}

\section{Założenia architektoniczne i izolacja sieciowa}\label{sec:izolacja_sieciowa}

Fundamentem eksperymentu jest utworzenie odizolowanego, powtarzalnego środowiska wirtualnego, które uniemożliwia wyciek szkodliwego kodu poza kontur badawczy, a jednocześnie dostarcza analizowanym próbkom wiarygodne syntetyczne odpowiedzi sieciowe. Schemat topologii oraz podział ról w sieci przedstawiono na rysunku~\ref{fig:topologia_lab}.

\Figure[0.85]{topology.png}{Topologia połączeń i podział ról w odizolowanej podsieci laboratoryjnej}{fig:topologia_lab}

Do realizacji projektu wykorzystano hiperwizor typu 2 (\textit{Type-2 Hypervisor}) VMware Workstation Pro. Wybór tego środowiska wynika z natywnego wsparcia dla formatu pamięci migawek (\file{.vmem}), który jest bezpośrednio odczytywany przez frameworki \textit{Volatility 3} oraz \textit{MemProcFS} bez konieczności uruchamiania dodatkowych agentów akwizycyjnych wewnątrz infekowanej maszyny.

Współpraca węzłów odbywa się w ramach dedykowanej sieci wirtualnej typu \textit{Host-Only} (\code{VMnet2}) z wyłączonym serwerem DHCP hiperwizora oraz całkowicie zablokowanym routingiem do fizycznej karty sieciowej hosta. Adresacja IP została skonfigurowana statycznie i ujęta w tabeli~\ref{tab:adresacja_sieciowa}.

\begin{table}[htbp]
\centering
\caption{Konfiguracja adresacji wirtualnej podsieci badawczej}
\label{tab:adresacja_sieciowa}
\tablefont
\begin{tabularx}{\linewidth}{l l L{3.5cm} Y}
\toprule
\hcell{Węzeł} & \hcell{Adres IP} & \hcell{Rola w sieci} & \hcell{Opis interfejsu} \\
\midrule
REMnux & \code{192.168.100.1} & Brama domyślna / DNS & Pętla symulacji C2 (INetSim / FakeNet-NG) \\
Windows Victim & \code{192.168.100.10} & Stacja ofiary & Punkt detonacji próbek złośliwego oprogramowania \\
\bottomrule
\end{tabularx}
\end{table}

\section{Przygotowanie stacji ofiary i konfiguracja telemetrii}\label{sec:konfiguracja_windows}

Maszyna ofiary bazuje na 64-bitowym systemie Microsoft Windows 10/11 Pro. Celem zapewnienia niezakłóconego wykonania próbek złośliwego oprogramowania oraz pełnej rejestracji artefaktów podjęto opisane poniżej kroki konfiguracyjne.

\subsection{Specyfikacja systemu i dezaktywacja mechanizmów obronnych}\label{subsec:dezaktywacja_obrony}

Aplikacje typu \textit{Infostealer} oraz \textit{RAT} często przerywają działanie w przypadku wykrycia aktywnych mechanizmów ochronnych. Aby zagwarantować pełne wykonanie ładunku (stadium \textit{post-exploitation}), wyłączono następujące moduły:
\begin{enumerate}
    \item \textbf{Windows Defender i SmartScreen:} Trwale zdezaktywowane (po uprzednim wyłączeniu mechanizmu \textit{Tamper Protection} w interfejsie zabezpieczeń Windows) za pomocą lokalnych polis grup (GPO) oraz modyfikacji kluczy rejestru przedstawionych na listingu~\ref{lst:disable_defender}.
    \item \textbf{Antimalware Scan Interface (AMSI):} Wyłączone skanowanie skryptów w czasie rzeczywistym, aby umożliwić pełne wykonanie obfuskowanych loaderów PowerShell oraz VBScript/JScript.
    \item \textbf{Windows Update oraz Windows Firewall:} Zablokowane powiązane usługi systemowe w celu uniknięcia niekontrolowanych zmian w strukturach jądra i wersjach bibliotek DLL oraz umożliwienia swobodnej komunikacji wychodzącej do lokalnej bramy C2.
\end{enumerate}

\begin{lstlisting}[style=terminal, caption={Dezaktywacja ochrony Windows Defender i zapory z poziomu PowerShell}, label={lst:disable_defender}]
Set-MpPreference -DisableRealtimeMonitoring $true -DisableIOAVProtection $true -DisableBehaviorMonitoring $true -DisableBlockAtFirstSeen $true -DisableScriptScanning $true
New-ItemProperty -Path "HKLM:\SOFTWARE\Policies\Microsoft\Windows Defender" -Name "DisableAntiSpyware" -Value 1 -PropertyType DWORD -Force
Set-NetFirewallProfile -Profile Domain,Public,Private -Enabled False
\end{lstlisting}

\subsection{Wdrożenie i profilowanie rozszerzonej telemetrii Sysmon}\label{subsec:wdrozenie_sysmon}

Standardowe dzienniki Windows Event Log nie dostarczają dostatecznej szczegółowości w obszarze tworzenia zdalnych wątków czy dostępu do pamięci procesów. Na stacji wdrożono narzędzie \textbf{Microsoft Sysmon} z profilowaniem opartym na standardzie \textit{Sysmon-Modular} (listing~\ref{lst:sysmon_install}).

\begin{lstlisting}[style=terminal, caption={Instalacja usługi Sysmon z dedykowanym plikiem konfiguracyjnym}, label={lst:sysmon_install}]
Sysmon64.exe -accepteula -i sysmonconfig-export.xml
\end{lstlisting}

Kluczowe identyfikatory zdarzeń (Event ID) skonfigurowane w profilu pod kątem analizy poinfekcyjnej ujęto w tabeli~\ref{tab:sysmon_events}.

\begin{table}[htbp]
\centering
\caption{Konfiguracja reguł audytu Sysmon na stacji badawczej}
\label{tab:sysmon_events}
\tablefont
\begin{tabularx}{\linewidth}{c L{4.5cm} Y}
\toprule
\hcell{Event ID} & \hcell{Nazwa zdarzenia} & \hcell{Cel analityczny / wykrywane techniki} \\
\midrule
1 & ProcessCreate & Śledzenie drzewa procesów, parametrów CLI, wywołań LOLBins. \\
3 & NetworkConnect & Powiązanie połączeń TCP/UDP z konkretnym identyfikatorem procesu (\code{PID}). \\
6 & DriverLoad & Wykrywanie ładowania nielicencjonowanych sterowników (rootkity / BYOVD). \\
7 & ImageLoad & Rejestracja ładowanych bibliotek DLL (wsparcie detekcji iniekcji i DLL Hijacking). \\
8 & CreateRemoteThread & Kluczowe dla detekcji \textit{Process Injection} oraz \textit{Process Hollowing}. \\
10 & ProcessAccess & Wykrywanie odczytu pamięci \file{lsass.exe} oraz manipulacji pamięcią obcych procesów. \\
11 & FileCreate & Rejestracja zrzucanych ładunków wykonywalnych oraz plików persystencji. \\
12/13 & RegistryEvent & Monitorowanie kluczy autostartu (\code{Run}, \code{RunOnce}, \code{Services}). \\
22 & DNSEvent & Rejestracja zapytań DNS generowanych przez procesy złośliwego oprogramowania. \\
\bottomrule
\end{tabularx}
\end{table}

\subsection{Konfiguracja natywnych dzienników Windows i buforów EVTX}\label{subsec:optymalizacja_evtx}

Standardowy rozmiar pliku dziennika \file{Microsoft-Windows-Sysmon/Operational.evtx} (20~MB) jest niewystarczający podczas intensywnej infekcji i może prowadzić do nadpisywania zdarzeń. Zwiększono limity buforów do 256~MB za pomocą narzędzia \cmd{wevtutil} (listing~\ref{lst:wevtutil_config}) oraz aktywowano pełne logowanie bloków skryptowych PowerShell (\textit{Script Block Logging}, Event ID 4104).

\begin{lstlisting}[style=terminal, caption={Zwiększenie maksymalnego rozmiaru buforów dzienników zdarzeń EVTX}, label={lst:wevtutil_config}]
wevtutil sl "Microsoft-Windows-Sysmon/Operational" /ms:268435456
wevtutil sl "Security" /ms:268435456
wevtutil sl "Microsoft-Windows-PowerShell/Operational" /ms:268435456
\end{lstlisting}

\section{Konfiguracja węzła analityczno-symulacyjnego (REMnux)}\label{sec:konfiguracja_remnux}

Węzeł analityczny oparto na dystrybucji REMnux (system Linux bazujący na Ubuntu), stanowiącej wyspecjalizowane środowisko do inżynierii wstecznej, analizy złośliwego oprogramowania i informatyki śledczej.

\subsection{Emulacja usług sieciowych i przechwytywanie ruchu (INetSim / FakeNet-NG, tshark)}\label{subsec:emulacja_sieci}

Większość współczesnych próbek złośliwego oprogramowania przed przejściem do właściwej fazy ataku weryfikuje dostępność łączności sieciowej poprzez zapytania DNS lub pobranie testowego zasobu HTTP. Brak odpowiedzi powoduje przedwczesne zakończenie procesu.

Na stacji REMnux skonfigurowano usługę \textbf{INetSim} (wspieraną w razie potrzeby przez \textbf{FakeNet-NG}), która nasłuchuje na interfejsie \code{192.168.100.1} i emuluje standardowe protokoły aplikacyjne:
\begin{itemize}
    \item \textbf{Serwer DNS:} Odpowiada adresem bramy (\code{192.168.100.1}) na wszystkie zapytania o nazwy domenowe kierowane przez malware do serwerów C2.
    \item \textbf{Serwer HTTP/HTTPS:} Zwraca poprawne kody odpowiedzi (\code{200 OK}) oraz serwuje syntetyczne pliki w odpowiedzi na żądania pobrania ładunków kolejnego etapu (\textit{second-stage payloads}).
    \item \textbf{Serwery SMTP / FTP / Dummy TCP:} Przechwytują próby eksfiltracji wykradzionych danych oraz surowe sesje komunikacyjne RAT.
\end{itemize}

Cały ruch sieciowy na interfejsie wewnętrznym jest równolegle rejestrowany do plików \file{.pcap} za pomocą narzędzia \cmd{tshark} (listing~\ref{lst:tshark_capture}).

\begin{lstlisting}[style=terminal, caption={Uruchomienie pełnego przechwytywania pakietów na interfejsie laboratoryjnym}, label={lst:tshark_capture}]
tshark -i eth0 -w /analiza/dumps/capture_infekcja.pcap
\end{lstlisting}

\subsection{Stos narzędziowy analizy śledczej i inżynierii detekcji}\label{subsec:stos_narzedziowy}

Na maszynie REMnux oraz hoście analitycznym przygotowano zestaw oprogramowania podzielony na cztery obszary funkcjonalne:
\begin{enumerate}
    \item \textbf{Wstępna analiza statyczna (Triage):} \textit{PEstudio}, \textit{Detect It Easy (DIE)}, \textit{FLOSS} (ekstrakcja obfuskowanych ciągów znaków) oraz deasembler/dekompilator \textit{Ghidra}.
    \item \textbf{Analiza pamięci operacyjnej RAM:} \textit{Volatility 3} (wraz z tablicami symboli PDB dla docelowego jądra Windows) oraz \textit{MemProcFS} umożliwiający zamontowanie zrzutu pamięci jako wirtualnego systemu plików.
    \item \textbf{Triage i korelacja dzienników zdarzeń:} \textit{Hayabusa} oraz \textit{Chainsaw} wyposażone w bazę reguł detekcyjnych \textit{SigmaHQ}.
    \item \textbf{Inżynieria i walidacja reguł:} Silnik \textit{YARA 4.x} oraz narzędzie \cmd{sigma-cli} do konwersji i weryfikacji składniowej reguł Sigma.
\end{enumerate}

\section{Metodyka prowadzenia eksperymentu i akwizycji artefaktów}\label{sec:metodyka_akwizycji}

W celu zachowania rygoru badawczego i pełnej porównywalności wyników, dla każdej badanej próbki złośliwego oprogramowania (opisanej w rozdziałach~\ref{ch:studium_przypadku_1} oraz~\ref{ch:studium_przypadku_2}) stosowana jest jednolita procedura eksperymentalna.

\subsection{Kryteria doboru próbek i procedura wstępnej weryfikacji statycznej}\label{subsec:dobor_probek}

Materiał badawczy pozyskiwany jest z publicznego repozytorium próbek złośliwego oprogramowania \textit{MalwareBazaar} (\code{abuse.ch}). Dobór obejmuje aktywne próbki z rodzin koni trojańskich zdalnego dostępu (RAT), loaderów oraz programów wykradających dane (Infostealer) przeznaczonych dla architektury x64/x86 systemu Windows. Przed uruchomieniem każdej próbki wykonywana jest analiza statyczna obejmująca:
\begin{itemize}
    \item wyznaczenie kryptograficznych skrótów identyfikacyjnych (MD5, SHA-1, SHA-256),
    \item inspekcję nagłówków PE, sekcji wykonywalnych, współczynnika entropii (weryfikacja spakowania/obfuskacji) oraz tablicy importów (IAT),
    \item ekstrakcję jawnych oraz dekodowanych ciągów znaków za pomocą narzędzia \textit{FLOSS}.
\end{itemize}

\subsection{Procedura kontrolowanej detonacji i akwizycji poinfekcyjnej}\label{subsec:procedura_detonacji}

Każdy eksperyment realizowany jest w sześciu powtarzalnych etapach zestawionych w tabeli~\ref{tab:cykl_badawczy}. Zrzut pamięci operacyjnej (\file{.vmem}) pozyskiwany jest poprzez wstrzymanie maszyny wirtualnej (\textit{Suspend}) lub wykonanie migawki z pamięcią na poziomie hiperwizora VMware, co eliminuje ryzyko wykrycia narzędzia akwizycyjnego przez złośliwy kod (\textit{anti-forensics}) i gwarantuje atomową spójność obrazu RAM.

\begin{table}[htbp]
\centering
\caption{Iteracyjny cykl prowadzenia eksperymentu badawczego i akwizycji materiału dowodowego}
\label{tab:cykl_badawczy}
\tablefont
\begin{tabularx}{\linewidth}{c L{4.2cm} Y}
\toprule
\hcell{Etap} & \hcell{Faza procedury} & \hcell{Wykonywane czynności operacyjne} \\
\midrule
1 & Przywrócenie stanu bazowego & Odtworzenie zweryfikowanej migawki \code{Clean_State} na stacji Windows Victim oraz wyczyszczenie logów usług na węźle REMnux. \\
\addlinespace
2 & Inicjalizacja nasłuchu & Uruchomienie emulacji usług w pakiecie INetSim oraz rejestracji pakietów \cmd{tshark} na interfejsie bramy analitycznej. \\
\addlinespace
3 & Kontrolowana detonacja & Przeniesienie próbki na stację ofiary i jej uruchomienie z uprawnieniami użytkownika lub administratora. \\
\addlinespace
4 & Okno obserwacji poinfekcyjnej & Utrzymanie działania systemu przez 5-10 minut w celu umożliwienia rozpakowania ładunku w RAM, iniekcji procesowej, założenia persystencji i nawiązania sesji C2. \\
\addlinespace
5 & Akwizycja artefaktów & Eksport dzienników \file{.evtx} (Sysmon, Security, PowerShell), zatrzymanie przechwytywania \file{.pcap} oraz wykonanie pełnego zrzutu pamięci RAM (\file{.vmem}) z poziomu hiperwizora. \\
\addlinespace
6 & Weryfikacja integralności & Obliczenie sum kontrolnych SHA-256 dla wszystkich zebranych plików dowodowych i przekazanie ich do analizy na węźle REMnux. \\
\bottomrule
\end{tabularx}
\end{table}

\subsection{Metodyka korelacji wieloźródłowej oraz walidacji reguł detekcyjnych (FP/TP)}\label{subsec:metodyka_walidacji}

Zebrany materiał dowodowy poddawany jest dwutorowej analizie śledczej:
\begin{enumerate}
    \item \textbf{Ścieżka analizy pamięci RAM:} Identyfikacja anomalii w strukturach procesów (\code{pslist}, \code{pstree}, \code{psscan}, \code{cmdline}), lokalizacja wstrzykniętego kodu w obszarach pamięci prywatnej z uprawnieniami \code{PAGE_EXECUTE_READWRITE} (\code{malfind}, \code{vadinfo}, MemProcFS) oraz ekstrakcja odszyfrowanych obrazów PE bezpośrednio z przestrzeni adresowej procesu.
    \item \textbf{Ścieżka analizy telemetrii i logów:} Automatyczny triage plików \file{.evtx} narzędziami \textit{Hayabusa} i \textit{Chainsaw}, a następnie szczegółowa rekonstrukcja łańcucha zdarzeń Sysmon i korelacja z ruchem sieciowym \file{.pcap}.
\end{enumerate}

Integracja obu ścieżek następuje poprzez powiązanie identyfikatorów procesów (\code{PID}, \code{ParentPID}), unikalnych deskryptorów \code{ProcessGuid}, znaczników czasu (\textit{timestamps} UTC) oraz gniazd sieciowych. Na tej podstawie wytwarzane są autorskie reguły **YARA** (ukierunkowane na artefakty binarne i ciągi znaków wyekstrahowane z pamięci RAM) oraz reguły **Sigma** (opisujące wzorce behawioralne w dziennikach zdarzeń).

Każda wytworzona reguła przechodzi dwuetapową walidację empiryczną:
\begin{itemize}
    \item \textbf{Test skuteczności (\textit{True Positive}, TP):} Potwierdzenie poprawnego wyzwolenia reguły na poinfekcyjnym zrzucie pamięci RAM (\file{.vmem}) lub zbiorze dzienników \file{.evtx} pochodzących z zainfekowanego systemu.
    \item \textbf{Test odporności na fałszywe alarmy (\textit{False Positive}, FP):} Przeskanowanie czystego zrzutu pamięci referencyjnej oraz dzienników zdarzeń wygenerowanych podczas standardowej pracy systemu (migawka bazowa \code{Clean_State}) w celu potwierdzenia zerowej liczby fałszywych trafień.
\end{itemize}

